% Solutions to the exercises of Chapter 27.
\begin{solutions}
  \item At a random clone pose and a feature with $\hat z = 0.3$, central differences with $\varepsilon = 10^{-6}$ agree with \eqref{eq:planar-H} to $10^{-10}$ in all five entries. For the
        heading entry, $\beta = \arctan z$ is the angle between the body axis $x_b$ and the ray to the feature. Turning the robot left by $\delta\psi$ turns the body axis
        toward the ray, so every bearing \emph{decreases} by exactly $\delta\psi$, whatever the depth: $\delta\beta = -\delta\psi$. Since $z = \tan\beta$,
        \[
          \frac{\partial z}{\partial\delta\psi} = \frac{\dd z}{\dd\beta}\,\frac{\partial\beta}{\partial\delta\psi} = \frac{1}{\cos^2\beta}\,(-1) = -(1 + \tan^2\beta) = -(1 + z^2).
        \]
        The factor $1 + z^2$ only converts angle to normalized coordinate: at the edge of the $\pm50^\circ$ field of view it is $2.4$, while in angle every point moves
        by the same $\delta\psi$.

  \item The top-right block of $I + F_c\Dt$ is $\Omega R(\hat\psi_k)e_1v_m\Dt$, while that of $\Phi_k$ is $\Omega(\hat p_{k+1} - \hat p_k) = \Omega R(\hat\psi_k + \tfrac12\omega_m\Dt)e_1v_m\Dt$. With
        $R(\hat\psi_k + \tfrac12\omega_m\Dt) = R(\hat\psi_k)(I + \tfrac12\omega_m\Dt\,\Omega) + O(\Dt^2)$ and $\Omega^2 = -I$,
        \[
          \Phi_k - (I + F_c\Dt) = \begin{bmatrix} 0 & -\tfrac12v_m\omega_m\Dt^2\,R(\hat\psi_k)e_1 \\ 0 & 0\end{bmatrix} + O(\Dt^3).
        \]
        The term is $\tfrac12\Dt^2\,\Omega\ddot{\hat p}$: the nominal path turns during the step, so its displacement is not $\dot{\hat p}_k\Dt$. Note that $F_c^2 = 0$, so for a
        \emph{constant} $F_c$ the matrix exponential would add nothing. The whole difference comes from $F_c$ changing during the step: $I + F_c\Dt$ is the rectangle rule for
        $I + \int F_c\,\dd t$. For the simulation ($v = 1$~m/s, $\omega = 0.1$~rad/s, $\Dt = 0.01$~s) it is $5\times10^{-6}$ per step, relative size $\tfrac12\omega\Dt = 5\times10^{-4}$. It is small,
        but it means the approximate transitions no longer telescope exactly into $\Omega(\hat p_k - \hat p_0)$ (\cref{ch:observability}).

  \item With $\omega = \omega_m - b_\omega - n_\omega$ and the nominal $\dot{\hat\psi} = \omega_m - \hat b_\omega$, $\dot{\delta\psi} = -\delta b_\omega - n_\omega$ and $\dot{\delta b}_\omega = n_b$; the position row is unchanged. For
        $\delta x_R = (\delta p, \delta\psi, \delta b_\omega)$ and $n = (n_v, n_\omega, n_b)$,
        \[
          F_c = \begin{bmatrix} 0_{2\times2} & \Omega R(\hat\psi)e_1v_m & 0_{2\times1} \\ 0 & 0 & -1 \\ 0 & 0 & 0 \end{bmatrix}, \qquad
          G_c = \begin{bmatrix} -R(\hat\psi)e_1 & 0 & 0 \\ 0 & -1 & 0 \\ 0 & 0 & 1 \end{bmatrix},
        \]
        with $Q_c = \diag(\sigma_v^2, \sigma_\omega^2, \sigma_b^2)$: the new column of $F_c$ is $(0, 0, -1, 0)\T$ and the new row of $G_c$ is $(0, 0, 1)$. The camera measures relative headings of
        the clones directly: the heading column $-(1 + z^2)$ does not depend on depth, and turning a single clone is visible. A bias error makes the heading error grow
        linearly across the window, $\delta\psi(t) = \delta\psi_0 - t\,\delta b_\omega$. The unobservable directions cannot produce such a ramp: a global rotation adds the same constant to every
        heading, and translations and scale leave headings unchanged. So the ramp changes the projected residuals, and the bias is observable. Numerically: on
        the simulated circle with ten clones and six landmarks, the observability matrix of $(\delta p, \delta\psi, \delta b_\omega)$ at the first clone time has rank one. Its
        three-dimensional nullspace (the two translations and the rotation) has bias components of $10^{-14}$.
        \ifsolutions
        In three dimensions the attitude row is $\dot{\delta\theta} = -\sx{\hat\omega}\delta\theta - \delta b_g - n_g$ and the bias row $\dot{\delta b}_g = n_{b_g}$. The $-I$ in the gyro-bias columns is the exact
        counterpart of the $-1$ here. The extra term $-\sx{\hat\omega}\delta\theta$ has no planar analog, because a planar heading error does not rotate with the body. For the same
        reason as here, all three gyroscope-bias components are observable once the camera measures relative attitudes: even the bias about the vertical,
        although the global yaw is not.
        \else
        In three dimensions the same mechanism makes all three gyroscope-bias components observable once the camera measures relative attitudes: even the bias
        about the vertical, although the global yaw is not. (The corresponding block of $F_c$ is part of the derivation exercises of \cref{ch:errordyn}.)
        \fi

  \item The body-frame points are $(4, 1)$, $(4, 0)$, $(4, -2)$, so $\hat z = (\tfrac14, 0, -\tfrac12)$ and
        \[
          H_f = \begin{bmatrix} -\tfrac1{16} & \tfrac14 \\ 0 & \tfrac14 \\ \tfrac18 & \tfrac14 \end{bmatrix}, \qquad
          h_1\times h_2 = \tfrac1{64}(-2, 3, -1), \qquad a = \tfrac1{\sqrt{14}}(2, -3, 1),
        \]
        so $r_o = (2r_1 - 3r_2 + r_3)/\sqrt{14}$. For sideways motion at a fixed heading, $z_i = (y_f - y_i)/x_f$ is a linear function of the camera position $y_i = (-1, 0, 2)$, and
        the weights annihilate exactly the constants and the linear functions: $\sum_ia_i = 0$ and $\sum_ia_iy_i = 0$. For baselines $b_1 = y_2 - y_1$ and $b_2 = y_3 - y_2$, the weights are
        $(b_2, -(b_1 + b_2), b_1)$, and $r_o = 0$ exactly when
        \[
          \frac{z_1 - z_2}{z_2 - z_3} = \frac{b_1}{b_2} .
        \]
        The depth $x_f$ cancels in the ratio of image displacements, so the residual compares the observed ratio with the estimated ratio of the baselines, and
        nothing else. If the middle clone were really at $y = 0.1$, the observed ratio would be $1.1/1.9 = 0.58$ instead of $\tfrac12$, and $r_o = -0.020$. The projected Jacobian
        is $H_o = \tfrac1{\sqrt{14}}\big[(\tfrac18, -\tfrac12, -\tfrac{17}8)\ \ (0, \tfrac34, 3)\ \ (-\tfrac18, -\tfrac14, -\tfrac54)\big]$; translation, rotation and scale of the configuration give zero, as before.

  \item With $a = (1, -2, 1)/\sqrt6$, the $\delta p_x$ entries of $H_o$ are $\tfrac{1}{16\sqrt6}$, $0$ and $-\tfrac{1}{16\sqrt6}$. Moving the middle clone along $x_w$ is invisible; moving the first clone
        along $x_w$ is visible, though weakly. The middle clone looks straight at the feature ($\hat z_2 = 0$), so $x_w$ is its own viewing ray, the nullspace of its projection
        Jacobian: moving along the ray changes the depth, which one view cannot see, and leaves the bearing unchanged. The first clone sees the feature at
        $\hat z_1 = \tfrac14$ ($14^\circ$). Moving it along $x_w$ changes that bearing by $\tfrac{1}{16}$ per meter, and the feature cannot absorb the change without disturbing the other two bearings. In
        general, moving a clone along its ray to a feature is invisible in that feature's residual. Other directions are visible, with a sensitivity
        proportional to the sine of the angle between motion and ray, divided by the depth.

  \item Each observation gives as many numbers as the image has dimensions, and the feature has as many unknowns as the space. A 1D camera in the plane: $H_{f,j}$ is
        $M\times2$ of rank 2, leaving $M - 2$ constraints, so $M = 2$ gives none. Two rays in a plane always intersect (unless parallel), so any two bearings are consistent with some
        feature. A 2D camera in space: $H_{f,j}$ is $2M\times3$ of rank 3, leaving $2M - 3$, so $M = 2$ gives one. Two rays in space are generically skew, and requiring them to intersect,
        i.e.\ to be coplanar with the baseline, is the epipolar constraint (Exercise~26.3).

  \item Final RMS errors after 90~s over 30 runs (dead reckoning: $2.80$~m):
        \begin{center}\small
        \begin{tabular}{@{}l c c c c@{}}
        \toprule
        $N_{\max}$ (window length) & 3 (0.3~s) & 5 (0.5~s) & 10 (1~s) & 20 (2~s) \\
        \midrule
        position & 2.23~m & 2.13~m & 0.85~m & 0.33~m \\
        heading & 0.19~rad & 0.17~rad & 0.069~rad & 0.025~rad \\
        average NEES (last 9~s, 3 DOF) & 2.65 & 2.97 & 3.14 & 3.20 \\
        \bottomrule
        \end{tabular}
        \end{center}
        The position error follows the heading error, which turns into cross-track error. The key quantity is how well one window determines the heading
        change across it, a heading ramp, which the camera sees only through parallax. With the scenario's landmarks, a single window determines it with a
        standard deviation of $0.05$, $0.022$, $0.008$ and $0.003$~rad for $N_{\max} = 3$, 5, 10, 20. Odometry alone has $\sigma_\omega\sqrt T = 0.009$, $0.013$, $0.019$ and $0.028$~rad over the same spans.
        Short windows are worse than odometry at exactly the quantity that drives the drift; the camera wins only for windows longer than about half a second,
        which is where the error drops. The benefit then levels off: $0.18$~m at $N_{\max} = 40$ and $0.17$~m at 80. Several things limit it:
        \begin{itemize}
          \item once the window pins the heading change far better than odometry, the remaining error is the unobservable drift of global heading and
                position, which no window removes;
          \item tracks end (here a landmark stays in view for a median of 6.8~s), so a longer window cannot lengthen them;
          \item the cost grows as $n^3$ with $n = 3 + 3N_{\max}$;
          \item old clones are never relinearized;
          \item the NEES creeps upward ($3.24$ at $N_{\max} = 80$), as more constraints feed more spurious heading information (\cref{ch:observability}).
        \end{itemize}

  \item Let $h(x, p_f) = \pi\big(R(\psi_i)\T(p_f - p_i)\big)$, or any function of $R(\psi_i)\T(p_f - p_i)$. Translating everything by $t$ ($p_i \to p_i + t$, $p_f \to p_f + t$) leaves $p_f - p_i$ unchanged.
        Rotating everything about the origin by $\delta$ ($p_i \to R(\delta)p_i$, $\psi_i \to \psi_i + \delta$, $p_f \to R(\delta)p_f$) gives
        $R(\psi_i + \delta)\T R(\delta)(p_f - p_i) = R(\psi_i)\T R(\delta)\T R(\delta)(p_f - p_i) = R(\psi_i)\T(p_f - p_i)$. So $h$ is constant along both families, identically in $t$ and $\delta$.
        Differentiating at $t = 0$ and $\delta = 0$, at the common linearization point $(\hat x, \hat p_f)$: the state moves along $N_t$ and the feature by $I$, or the state along $N_r$ and the
        feature by $\Omega\hat p_f$. Hence
        \[
          H_{x,j}N_t + H_{f,j} = 0, \qquad H_{x,j}N_r + H_{f,j}\,\Omega\hat p_f = 0,
        \]
        and multiplying by $A_j\T$ gives $H_{o,j}N = 0$. Nothing about $\pi$ was used. The proof also shows why the hypothesis matters: the identities come from differentiating one
        function at one point, so they hold only if every block of $H_{x,j}$ and $H_{f,j}$ is evaluated at the same estimates.

  \item Each observation $(\|d_i\|, \atantwo(d_y, d_x))$ now gives two numbers for two unknowns of the feature, so a feature seen $M$ times gives $2M - 2$ constraints. A single
        observation determines the feature and constrains nothing; two observations already give two constraints. None of the directions of
        \eqref{eq:planar-nullspace} becomes observable: the measurement is still a function of $R(\psi_i)\T(p_f - p_i)$ alone, so the argument of Exercise~8 applies unchanged, and global
        translation and heading remain invisible. What changes is scale. The range measures distances, so the shape \emph{and size} of the trajectory are now
        observable from the camera alone, and the along-track error no longer depends on the odometry speed noise only.
\end{solutions}
